% Generated draft: requires mathematical and source review.
\documentclass[11pt]{article}
\usepackage[margin=25mm]{geometry}
\usepackage{fontspec}
\setmainfont{Latin Modern Roman}
\usepackage{amsmath,amssymb,mathtools,unicode-math}
\setmathfont{Latin Modern Math}
\usepackage{enumitem,needspace,xurl,hyperref}
\setlength{\parindent}{0pt}
\setlength{\parskip}{5pt}
\setlength{\emergencystretch}{3em}
\newcommand{\N}{\mathbb{N}}
\newcommand{\Z}{\mathbb{Z}}
\newcommand{\Q}{\mathbb{Q}}
\newcommand{\R}{\mathbb{R}}
\newcommand{\C}{\mathbb{C}}
\newcommand{\F}{\mathbb{F}}
\newcommand{\A}{\mathbb{A}}
\newcommand{\PP}{\mathbb{P}}
\newcommand{\E}{\mathbb{E}}
\newcommand{\eps}{\varepsilon}
\newcommand{\dd}{\,\mathrm{d}}
\begin{document}
\section*{UnsolvedMath ID: 30001575}
\textbf{Source problem code:} catalogue entry
\subsection{Definitions and mathematical statement}
For a complex algebraic surface \(S\), let \(\C(S)\) denote its field of rational functions. A smooth projective surface is minimal if it contains no smooth rational curve of self-intersection \(-1\).

The supplied record asks, ``What are Burniat surfaces?'' It then describes a Burniat surface as the minimal model associated with a specified nine-line configuration in \(\mathbb{P}^{2}\), with the following displayed function-field expression:
\[
\C\bigl(\sqrt{\Delta_1/\Delta_2},\sqrt{\Delta_1/\Delta_3}\bigr).
\]
This expression is retained as supplied; it does not specify a complete construction. The nine-line configuration and the meanings of \(\Delta_1,\Delta_2,\Delta_3\) are absent. In particular, it is unclear whether these symbols denote individual line equations or products of equations for groups of lines.

If the intended construction is a finite cover of the complex projective plane, its function field would need to be specified as an extension of \(\C(\mathbb{P}^{2})\). Conditional on \(\Delta_1/\Delta_2\) and \(\Delta_1/\Delta_3\) being well-defined nonzero rational functions on that plane, the corresponding candidate expression is
\[
\C(\mathbb{P}^{2})\bigl(\sqrt{\Delta_1/\Delta_2},\sqrt{\Delta_1/\Delta_3}\bigr).
\]
This is an explicitly conditional interpretation, not a correction verified against the report. The supplied verification describes the construction as a bidouble cover, but the missing configuration and equations prevent checking its degree and geometric hypotheses. No separate unresolved mathematical assertion or quantified conjecture is supplied.
\subsection{Short English statement}
The record gives an answered definitional question about Burniat surfaces, described through a nine-line configuration and a bidouble-cover construction. The configuration, the defining expressions, and the base function field are insufficiently specified to recover a precise construction. No distinct open problem can be formulated from the supplied text.
\subsection{Sources}
\begin{itemize}
\item Supplied local UnsolvedMath registry record, permanent ID \(30001575\), source code \texttt{catalogue entry}, title ``Answered Definition of Burniat Surfaces.'' No individual proposer is supplied; the recorded proposed year is 2010.
\item Supplied report citation: \emph{Komplexe Analysis (2011)}, Oberwolfach Reports, pp.~2283--2333, DOI \texttt{10.4171/owr/2010/39}. Supplied source URL and literature reference labelled ``original OWR report'': \url{https://doi.org/10.4171/owr/2010/39}.
\item Literature assessment dated exactly \texttt{2026-08-21}: the supplied curation classifies the record as open conservatively because the extracted text is a definitional question with no distinct claim to verify. This is supplied dated curation, not an independently verified literature assessment.
\item Supplied statement-verification note: the heading asks what Burniat surfaces are and the following text answers with a construction from nine-line configurations as bidouble covers. The record marks the statement \texttt{unrecoverable}, promotes no separate clean statement, and retains the original for provenance. The cited report has not been independently consulted for this draft.
\end{itemize}
\end{document}
